%%%%%%%%%%%%%%%%
% MARCO
%%%%%%%%%%%%%%%%

\section{Marco teórico y antecedentes}

\subsection{Bases Teóricas}

\begin{itemize}

    \item \textbf{Ingeniería de Software}  
    
    Enfoque sistemático y disciplinado para el diseño, desarrollo y mantenimiento de sistemas de software. A través de la aplicación de principios de ingeniería y buenas prácticas, esta disciplina busca construir soluciones confiables, eficientes y escalables que respondan a necesidades específicas de los usuarios o de una organización. Más allá de la implementación de código, la ingeniería de software abarca diversas actividades a lo largo del ciclo de vida de un sistema, incluyendo el análisis de requerimientos, el diseño de la arquitectura, la implementación, las pruebas y el mantenimiento continuo del software, con el fin de garantizar su correcto funcionamiento y evolución en el tiempo \citep{Jazayeri2003}.

    \item \textbf{Ciclo de vida del Software}  
    
    Proceso estructurado utilizado para planificar, desarrollar, probar, desplegar y mantener sistemas de software. Este enfoque organiza el desarrollo en una serie de fases que abarcan desde la conceptualización inicial de un producto hasta su implementación, operación, mantenimiento y eventual retiro. A través de esta estructura, el SDLC permite gestionar de manera ordenada las actividades involucradas en la construcción de software, asegurando que cada etapa produzca los resultados necesarios para la siguiente. De esta forma, el modelo facilita el control del proceso de desarrollo, mejora la eficiencia del trabajo de los equipos y contribuye a garantizar la calidad del software entregado \citep{Jazayeri2003}.

    \item \textbf{Domain Driven Design}  
    
    Enfoque que propone centraliza el desarrollo de software en el dominio del negocio, buscando que los modelos reflejen con precisión los procesos y objetivos críticos de la organización \citep{Vernon2013}. Este enfoque facilita la colaboración entre expertos de negocio y desarrolladores mediante un lenguaje común y promueve la estructuración de la aplicación en módulos y contextos delimitados con responsabilidades claras, lo que permite modularidad, desacoplamiento y reutilización de funcionalidades. Además, DDD orienta el diseño de la arquitectura y sus componentes para que estén alineados con el dominio, mejorando la mantenibilidad, la coherencia técnica y la capacidad de evolución controlada del sistema.

    \item \textbf{Modularidad en sistemas de Software}  
    
    Principio de desarrollar un sistema mediante la división en subsistemas más pequeños que pueden diseñarse e implementarse de manera relativamente independiente. Estos subsistemas pueden adoptar diferentes formas según el paradigma de desarrollo utilizado, como funciones en sistemas estructurados, clases en sistemas orientados a objetos, aspectos en programación orientada a aspectos o componentes en arquitecturas basadas en componentes. En términos generales, todos estos elementos pueden entenderse como unidades funcionales que encapsulan una parte específica de la lógica del sistema \citep{Yu2009}.

    \item \textbf{Arquitectura hexagonal}  
    
    Enfoque arquitectónico que propone organizar las aplicaciones de manera que la lógica central del sistema permanezca independiente de las tecnologías externas. 
    
    En este modelo, el núcleo de la aplicación contiene el dominio del sistema, es decir, las reglas de negocio y comportamientos que representan la funcionalidad principal de la aplicación. Al mantener este núcleo aislado de elementos como interfaces de usuario, bases de datos o servicios externos, se busca evitar que las decisiones tecnológicas afecten directamente la lógica del negocio \citep{AlistairCockburn2005}. 
    
    Para lograr esta separación, la arquitectura define puertos, que representan interfaces mediante las cuales el núcleo del sistema expone sus funcionalidades, y adaptadores, que son componentes encargados de conectar estas interfaces con tecnologías específicas. De esta manera, los adaptadores traducen las interacciones provenientes del exterior hacia las operaciones definidas por el dominio de la aplicación, permitiendo desacoplar la lógica de negocio de los detalles de infraestructura y facilitando la evolución y mantenimiento del sistema \citep{AlistairCockburn2005}.

    \item \textbf{Framework de desarrollo}  
    
    Conjunto estructurado de componentes de software reutilizables que provee una base y un conjunto de reglas para crear aplicaciones de forma más eficiente, consistente y estandarizada, facilitando el desarrollo mediante módulos predefinidos, APIs, librerías y herramientas que aceleran la construcción de funcionalidades sin necesidad de programar desde cero elementos comunes, como manejo de datos, control de flujo, seguridad o interfaces de usuario. Además, un framework ayuda a mantener la calidad, coherencia y consistencia del código, promueve buenas prácticas de desarrollo y facilita la organización y colaboración dentro de los equipos de desarrollo a lo largo del ciclo de vida de un proyecto \citep{aws-framework-2026}.

    \item \textbf{BIAN (Banking Industry Architecture Network)}  
    
    Asociación global independiente y sin fines de lucro formada por instituciones financieras, proveedores de software, consultoras, integradores y académicos con el propósito de definir y promover un marco arquitectónico común para la banca. La misión de BIAN es ofrecer un modelo estandarizado de referencia para la interoperabilidad de sistemas bancarios, apoyando la innovación tecnológica, reduciendo los costos de integración y simplificando la arquitectura de TI en el sector financiero. Para lograr esto, BIAN reúne experiencia de su comunidad para crear un marco de mejores prácticas que los bancos pueden utilizar al diseñar, implementar y probar servicios y plataformas financieras \citep{bian-2026}.
    
    \item \textbf{ISO 20022}  
    
    Estándar internacional de mensajería financiera desarrollado por el Technical Committee ISO/TC68 Financial Services de la Organización Internacional de Normalización (ISO). Este estándar proporciona una plataforma común para el desarrollo de mensajes financieros mediante un conjunto de reglas, un diccionario de datos y un catálogo de modelos que describen los procesos de negocio y los elementos de datos utilizados en las comunicaciones financieras. Los modelos y los mensajes derivados se publican en el ISO 20022 Financial Repository, lo que permite a las instituciones financieras definir conjuntos de mensajes basados en un enfoque internacionalmente acordado y consistente en cuanto a semántica y estructura \citep{iso20022-2026}.

\end{itemize}

\subsection{Antecedentes}

En el desarrollo de software, se han implementado diversas estrategias y herramientas para agilizar el inicio de proyectos y garantizar consistencia en las prácticas de desarrollo. Estas soluciones pueden organizarse jerárquicamente según su alcance y complejidad, desde los fundamentos conceptuales hasta plataformas completas de gestión del ciclo de vida de aplicaciones.


\subsubsection*{Fundamentos conceptuales}

\begin{itemize}

    \item \textbf{Scaffolding}  
    
    Técnica que permite generar automáticamente la estructura básica de una aplicación o de determinados componentes dentro de un proyecto. A partir de configuraciones o modelos definidos, las herramientas de scaffolding crean archivos, directorios y código inicial necesarios para implementar funcionalidades comunes. Esto permite a los desarrolladores contar con una base funcional sobre la cual continuar el desarrollo, acelerando las primeras etapas de construcción del sistema.

    Diversos frameworks modernos incorporan herramientas de scaffolding que permiten generar automáticamente módulos o componentes predefinidos, reduciendo el tiempo necesario para implementar nuevas funcionalidades. Por ejemplo, el framework NestJS proporciona una herramienta de línea de comandos (CLI) que permite generar elementos como módulos, controladores y servicios mediante comandos específicos, simplificando la creación de estructuras estándar dentro de las aplicaciones \citep{nestjs-2026}.

    \item \textbf{Boilerplate}  
    
    Fragmentos de código que pueden reutilizarse en diferentes contextos con pocas o ninguna modificación. Este término proviene originalmente del uso de plantillas en documentos, donde se reutilizan textos estándar que luego se adaptan ligeramente a situaciones específicas. De forma similar, en el desarrollo de software los programadores utilizan secciones de código previamente escritas para evitar tener que implementarlas nuevamente desde cero \citep{aws-boilerplate-2026}.

    El uso de boilerplate permite mejorar la eficiencia del desarrollo, ya que facilita la reutilización de código probado y reduce el tiempo necesario para implementar funcionalidades comunes. Además, contribuye a mantener consistencia en la estructura del software y disminuir errores derivados de la repetición manual de código \citep{aws-boilerplate-2026}.

\end{itemize}

\subsubsection*{Herramientas especializadas de generación de código}

\begin{itemize}

    \item \textbf{Yeoman}  
    
    Herramienta de código abierto diseñada para facilitar la creación de nuevos proyectos mediante el uso de generadores. Estos generadores actúan como complementos que pueden ejecutarse desde la línea de comandos para generar automáticamente proyectos completos o partes específicas de una aplicación. 
    
    El objetivo de Yeoman es ayudar a los desarrolladores a iniciar proyectos rápidamente, proporcionando estructuras iniciales, configuraciones y herramientas que incorporan buenas prácticas de desarrollo \citep{yeoman-2026}.

    \item \textbf{Nx.js}  
    
    Plataforma de desarrollo de código abierto orientada a la gestión de repositorios monolíticos (monorepos) y a la automatización de tareas dentro de proyectos de software. Esta herramienta proporciona un sistema de construcción que permite ejecutar y coordinar tareas como compilación, pruebas y análisis de código de manera eficiente, especialmente en proyectos que contienen múltiples aplicaciones o librerías dentro de un mismo repositorio \citep{nx-2026}.

    Además, la herramienta puede generar aplicaciones, librerías y componentes mediante plugins y generadores, facilitando la estandarización de la arquitectura y la organización del código en proyectos de gran escala \citep{nx-2026}.

\end{itemize}

\subsubsection*{Plataformas de desarrollo internas (IDP)}

\begin{itemize}

    \item \textbf{Backstage}  
    
    Framework de código abierto para organizar las herramientas, servicios y documentación utilizados en proyectos de software mediante portales de desarrollo. Originalmente creado por Spotify para enfrentar la creciente complejidad de su infraestructura interna, Backstage permite a los equipos desarrollar y gestionar software de manera más eficiente al proporcionar un punto único de interacción con todos los recursos técnicos de la organización \citep{backstage-2026}.

    Su base es un catálogo de software centralizado que contiene metadatos sobre servicios, librerías, pipelines de datos y otros componentes, facilitando su descubrimiento y administración. Además, incluye plantillas de software que permiten generar nuevos proyectos de forma estandarizada según las mejores prácticas de la organización, así como un sistema de documentación técnica integrado (TechDocs) que sigue un enfoque de "docs like code" \citep{backstage-2026}.

    \item \textbf{Port}  
    
    Plataforma diseñada para unificar y organizar todas las herramientas, datos y flujos de trabajo que los equipos de ingeniería utilizan durante el ciclo de vida del software. Es una solución enfocada en mejorar la experiencia del desarrollador proporcionando un punto centralizado donde encontrar, lanzar, operar y automatizar tareas relacionadas con el desarrollo y la infraestructura de software \citep{port-platform-2026}.

    Integra varios componentes clave que ayudan a gestionar el ecosistema de desarrollo de forma coherente y eficiente. Entre ellos se encuentran un catálogo de software que unifica metadatos de servicios, entornos y otros activos técnicos, acciones preconfiguradas que estandarizan operaciones repetitivas, un orquestador de flujos de trabajo que automatiza procesos y scorecards para medir la calidad, seguridad y cumplimiento de los estándares del equipo \citep{port-platform-2026}.
    
\end{itemize}